\chapter[Równania Naviera--Stokesa]{Równania Naviera--Stokesa: rozkład Helmholtza, rozwiązania Leraya--Hopfa i struktura wirów}
\label{chap:pde28}

\section{Równanie nieściśliwego płynu i wybór przestrzeni}
Niech $d\in\{2,3\}$, $\nu>0$, $T>0$ i $\T^d=(\R/2\pi\mathbb Z)^d$.
Szukamy prędkości $u:(0,T)\times\T^d\to\R^d$ i ciśnienia
$p:(0,T)\times\T^d\to\R$ spełniających
\begin{equation}\label{eq:pde28-ns}
 \partial_tu+(u\cdot\nabla)u-\nu\Delta u+\nabla p=0,
 \qquad \diver u=0,
 \qquad u(0)=u_0.
\end{equation}
Współczynnik $\nu$ jest lepkością kinematyczną, $u$ jest \emph{polem
prędkości}, natomiast ciśnienie jest mnożnikiem wymuszającym warunek
nieściśliwości. Przyjęcie torusa pozwala wyprowadzić zasadnicze
argumenty bez dodatkowych założeń o regularności brzegu. Dla
pól wektorowych symbol $\nabla u$ oznacza macierz
$(\partial_j u_i)_{i,j}$, a $\diver(u\otimes u)$ jest wektorem
o współrzędnych $\sum_j\partial_j(u_i u_j)$.

Zakładamy $\int_{\T^d}u_0\,\mathrm d\lambda_d(x)=0$.
Dla gładkich rozwiązań średnia prędkości zachowuje się
w czasie: z $\diver u=0$ wynika
$(u\cdot\nabla)u=\diver(u\otimes u)$, a całki
z pochodnych periodycznych znikają. Całkując
równanie~\eqref{eq:pde28-ns} w przestrzeni, otrzymujemy
$\frac{\mathrm d}{\mathrm dt}\int_{\T^d}u\,\mathrm d\lambda_d=0$.
Dla rozwiązań dystrybucyjnych analogiczny rachunek
przeprowadzamy, testując równanie stałym wektorem
w przestrzeni pomnożonym przez gładką funkcję czasu
o zwartym nośniku.
Jeżeli średnia wynosi $m=(2\pi)^{-d}\int_{\T^d}u_0\,
\mathrm d\lambda_d$, to zmiana układu odniesienia
\[
 v(t,x)=u(t,x+tm)-m,\qquad
 \widetilde p(t,x)=p(t,x+tm)
\]
daje równanie tej samej postaci dla $v$ o średniej zero.
Istotnie, dodatkowy składnik $m\cdot\nabla u$
w pochodnej czasowej znosi się ze składnikiem
$-m\cdot\nabla u$ w członie konwekcyjnym;
laplasjan i dywergencja komutują z translacją.
Dla pól o średniej zero używamy przestrzeni
\[
 \begin{aligned}
 H&=\Bigl\{v\in L^2(\T^d;\R^d):\
 \diver v=0\text{ w }\mathcal D'(\T^d),\
 \int_{\T^d}v\,\mathrm d\lambda_d=0\Bigr\},\\
 V&=H\cap H^1(\T^d;\R^d).
 \end{aligned}
\]
Przestrzeń $H$ jest domknięta w $L^2$: średnia
jest ciągłym odwzorowaniem liniowym, a warunek
dystrybucyjny jest przecięciem jąder ciągłych
funkcjonałów $v\mapsto\int_{\T^d}v\cdot\nabla\psi\,
\mathrm d\lambda_d$, $\psi\in C^\infty(\T^d)$.
Analogicznie $V$ jest domknięta w $H^1$.
Przyjmujemy współczynniki Fouriera unormowane przez
$\widehat v(k)=(2\pi)^{-d}\int_{\T^d}v(x)e^{-ik\cdot x}
\,\mathrm d\lambda_d(x)$, natomiast normy $L^2$
obliczamy względem pełnej miary Lebesgue'a.
Z twierdzenia Parsevala~\ref{thm:pde6-parseval}
i zerowej średniej dla $v\in V$ mamy
\[
 \|v\|_2^2=(2\pi)^d\sum_{k\ne0}|\widehat v(k)|^2
 \le (2\pi)^d\sum_{k\ne0}|k|^2|\widehat v(k)|^2
 =\|\nabla v\|_2^2.
\]
Na $V$ przyjmujemy iloczyn skalarny
$(v,w)_V=(\nabla v,\nabla w)_{L^2}$ i normę
$\|v\|_V=\|\nabla v\|_2$, równoważną normie $H^1$,
zatem $V$ jest przestrzenią Hilberta.
Skończone sumy Fouriera pól z $H$ zachowują zerową
średnią i warunek $k\cdot\widehat v(k)=0$,
a przy symetrycznym obcięciu częstości również
rzeczywistość pola. Są gęste w $H$, a dla $v\in V$
zbiegają do $v$ również w normie $V$.
Zatem zanurzenie $V\hookrightarrow H$ jest gęste
i ciągłe. Oznaczamy przez $V'$ przestrzeń dualną $V$
z normą $\|F\|_{V'}=\sup_{\|\varphi\|_V\le1}
|\langle F,\varphi\rangle|$.
Zanurzenie $H\hookrightarrow V'$ przyporządkowuje
$h$ funkcjonałowi $\varphi\mapsto(h,\varphi)_H$;
jest ciągłe i injektywne dzięki gęstości $V$.
Jego obraz jest gęsty w $V'$: gdyby nie był, istniałby
niezerowy ciągły funkcjonał na $V'$ zerujący się
na tym obrazie. Refleksywność $V$ pozwala zapisać go
jako ewaluację w pewnym $\varphi\in V\setminus\{0\}$,
co dawałoby $(h,\varphi)_H=0$ dla wszystkich $h\in H$.
Otrzymujemy trójkę przestrzeni Hilberta
$V\subset H\simeq H'\subset V'$ z gęstymi,
ciągłymi zanurzeniami.

\section{Rozkład Helmholtza i rzut Leraya}
\begin{twierdzenie}[Rozkład Helmholtza na torusie]\label{thm:pde28-helmholtz}
Niech $d\ge2$ (przestrzeń $H$ definiujemy jak powyżej
również dla tych wymiarów) i niech
$F\in L^2(\T^d;\R^d)$ ma średnią zero.
Istnieje dokładnie jedna para $(v,\phi)$, gdzie
$v\in H$, $\phi\in H^1(\T^d)$ ma średnią zero, taka że
$F=v+\nabla\phi$. Składniki są ortogonalne w $L^2$.
Przyporządkowanie $PF=v$ jest ortogonalnym rzutem
$L^2_0(\T^d;\R^d)\to H$, gdzie indeks $0$ oznacza
podprzestrzeń pól o średniej zero. Dla
$k\in\mathbb Z^d\setminus\{0\}$
\begin{equation}\label{eq:pde28-P}
 \widehat{PF}(k)=\left(I-\frac{k\otimes k}{|k|^2}\right)\widehat F(k),
 \qquad \widehat\phi(k)=-\frac{i k\cdot\widehat F(k)}{|k|^2}.
\end{equation}
\end{twierdzenie}
\begin{proof}
Przyjmujemy konwencję $\widehat F(k)=(2\pi)^{-d}
\int_{\T^d}F(x)e^{-ik\cdot x}\,\mathrm d\lambda_d(x)$.
Tak jak dla szeregów Fouriera w rozdziale~\ref{chap:pde-fourier}
współczynniki są uśredniane; normy $L^2$ użyte przy oszacowaniach
energii płynu odnoszą się jednak do pełnej miary $\lambda_d$,
bez dzielenia przez $(2\pi)^d$.
Dla każdego $k\ne0$ macierz w~\eqref{eq:pde28-P} jest rzutem
ortogonalnym na podprzestrzeń prostopadłą do $k$.
Ustalamy $\widehat{PF}(0)=\widehat\phi(0)=0$.
Ponieważ norma operatorowa macierzy nie przekracza $1$,
twierdzenie Parsevala daje ciąg współczynników
o skończonej sumie kwadratów, a zupełność bazy Fouriera
wyznacza $PF\in L^2$. Dla rzeczywistego $F$
mamy $\widehat F(-k)=\overline{\widehat F(k)}$.
Macierz rzutu jest rzeczywista i parzysta względem $k$,
więc współczynniki $PF$ mają tę samą symetrię.
Współczynniki $\phi$ również spełniają
$\widehat\phi(-k)=\overline{\widehat\phi(k)}$.
Obie funkcje są więc rzeczywiste.
Mamy $k\cdot\widehat{PF}(k)=0$ dla każdego $k$,
co jest równoważne $\diver PF=0$ w sensie dystrybucji
(przez różniczkowanie szeregów w sensie dystrybucyjnym).
Średnia $PF$ jest zerowa.
Z drugiej strony $ik\widehat\phi(k)=\widehat F(k)-\widehat{PF}(k)$.
Nierówność $|k||\widehat\phi(k)|\le|\widehat F(k)|$
oraz $|k|\ge1$ dla $k\ne0$ dają
\[
 \sum_k(1+|k|^2)|\widehat\phi(k)|^2
 \le 2\sum_k|\widehat F(k)|^2<\infty.
\]
Stąd $\phi\in H^1$ i równość $F=PF+\nabla\phi$
wynika z równości wszystkich współczynników Fouriera.
Dla $v\in H$ i $\phi\in H^1$ wybieramy gładkie
periodyczne wielomiany trygonometryczne
$\phi_n\to\phi$ w $H^1$. Definicja dywergencji
dystrybucyjnej daje $\int_{\T^d}v\cdot\nabla\phi_n\,
\mathrm d\lambda_d=0$. Nierówność Cauchy'ego--Schwarza
pozwala przejść do granicy, więc
$\int_{\T^d}v\cdot\nabla\phi\,\mathrm d\lambda_d=0$.
Jest to ortogonalność składników. Jeżeli $F=v_1+\nabla\phi_1
=v_2+\nabla\phi_2$, to różnica obu rozkładów leży jednocześnie
w $H$ i w jego ortogonalnym dopełnieniu, zatem znika.
Warunek zerowej średniej wyznacza jednoznacznie $\phi$:
jeśli $\nabla(\phi_1-\phi_2)=0$, to nierówność
Poincarégo daje $\phi_1=\phi_2$.
Ponieważ $PF\in H$ oraz $F-PF\perp H$, jest to
ortogonalny rzut na $H$, a więc $P^2=P$ i $P^*=P$.
Wzór na współczynniki definiuje również $P$ na
periodycznych przestrzeniach $H^s$ o średniej zero
dla $s\in\R$; mnożnik ma normę co najwyżej $1$
dla każdego $k$, więc jest kontrakcją w normie
Fouriera z wagą $(1+|k|^2)^s$.
\end{proof}

\begin{uwaga}[Operator Stokesa]
Na $H$ określamy $A=-P\Delta$ z dziedziną
$D(A)=H^2(\T^d;\R^d)\cap H$. Ponieważ $P$ komutuje
z $\Delta$ na torusie, mamy $Av=-\Delta v$, a
$(Av,v)_H=\|\nabla v\|_2^2$. Dokładniej, dla każdej
pary częstości $\{k,-k\}$, $k\ne0$, wybieramy bazę
ortonormalną $k^\perp$ w $\R^d$ i mnożymy jej wektory
przez odpowiednio unormowane $\cos(k\cdot x)$
i $\sin(k\cdot x)$. Otrzymujemy rzeczywistą bazę
ortonormalną $(w_j)$ przestrzeni $H$, złożoną
z funkcji własnych $A$ o wartościach własnych
$\lambda_j=|k_j|^2$. Numerujemy je niemalejąco,
z krotnościami, więc $\lambda_j\ge1$ i
$\lambda_j\to\infty$.
Dziedzina $D(A)$ odpowiada warunkowi
$\sum_j\lambda_j^2|(v,w_j)_H|^2<\infty$,
równoważnemu warunkowi $H^2$ na polach o średniej zero.
Operator jest zatem unitarnie równoważny rzeczywistemu
operatorowi mnożenia przez ciąg $(\lambda_j)$
na jego maksymalnej dziedzinie w $\ell^2$.
Jest samosprzężony: jeśli ciąg należy do dziedziny
operatora sprzężonego, testowanie przez wektory
jednostkowe wymusza, by ciąg pomnożony przez
$(\lambda_j)$ należał do $\ell^2$, co jest dokładnie
warunkiem dziedziny $A$.
Mamy $A\ge I$ i $A^{-1}$ jest ograniczony na $H$.
Obcięcie mnożnika $\lambda_j^{-1}$ do pierwszych $m$
współrzędnych daje operator skończonego rzędu,
którego odległość od $A^{-1}$ w normie operatorowej
wynosi $\sup_{j>m}\lambda_j^{-1}\to0$.
Stąd $A^{-1}$ jest zwarty. Ponadto
$\|v\|_V^2=\sum_j\lambda_j|(v,w_j)_H|^2$,
więc
$\sum_{j>m}|(v,w_j)_H|^2\le
\lambda_{m+1}^{-1}\|v\|_V^2$.
Ten jednostajny szacunek ogonów i skończenie wymiarowe
projekcje dowodzą zwartości zanurzenia $V\hookrightarrow H$.
Rzut $P$ nie jest mnożeniem przez stałą macierz w przestrzeni
fizycznej, lecz nielokalnym operatorem Fouriera.
\end{uwaga}

\section{Tożsamość energii i słaba postać równania}
Dla $a,b,c\in V$ określamy formę trójliniową
\[
 b(a,b,c)=\int_{\T^d}(a\cdot\nabla)b\cdot c\,\mathrm d\lambda_d(x).
\]
W wymiarach $2$ i $3$ zanurzenie Sobolewa
$H^1(\T^d)\hookrightarrow L^4(\T^d)$ jest ciągłe.
Stąd nierówność Höldera z wykładnikami $4,2,4$ daje
\[
 |b(a,b,c)|\le\|a\|_4\|\nabla b\|_2\|c\|_4
 \le C\|a\|_V\|b\|_V\|c\|_V.
\]
Forma jest zatem dobrze określona i ciągła na $V^3$.
Dla gładkich pól periodycznych, gdy $\diver a=0$,
\[
 b(a,b,c)+b(a,c,b)
 =\int_{\T^d}a\cdot\nabla(b\cdot c)\,
 \mathrm d\lambda_d
 =-\int_{\T^d}(\diver a)(b\cdot c)\,\mathrm d\lambda_d=0.
\]
Aproksymujemy trzy pola ich skończonymi sumami Fouriera
w $V$ i przechodzimy do granicy dzięki ciągłości formy.
Otrzymujemy $b(a,b,c)=-b(a,c,b)$, w szczególności
$b(a,b,b)=0$. Oznacza to, że konwekcja nie wnosi
wkładu do zmiany energii kinetycznej; za dyssypację
odpowiada człon lepkościowy.
Dla gładkiego rozwiązania testowanie równania przez $u$
i periodyczne całkowanie przez części dają bowiem
\[
 \tfrac12\frac{\mathrm d}{\mathrm dt}\|u\|_2^2
 +\nu\|\nabla u\|_2^2=0,
\]
gdyż $\int_{\T^d}\nabla p\cdot u\,\mathrm d\lambda_d
=-\int_{\T^d}p\,\diver u\,\mathrm d\lambda_d=0$.
Testowanie przez stałe w czasie $\varphi\in V$
i całkowanie względem czasu prowadzi do słabej
postaci równania podanej niżej.

\begin{definicja}[Rozwiązanie Leraya--Hopfa na torusie]\label{def:pde28-leray}
Niech $u_0\in H$, $\nu>0$ i $T>0$. Pole $u$ nazywamy
rozwiązaniem Leraya--Hopfa zagadnienia~\eqref{eq:pde28-ns} na $[0,T]$,
jeżeli $u\in L^\infty(0,T;H)\cap L^2(0,T;V)$,
$u\in C([0,T];H_{\mathrm w})$, $u(0)=u_0$ oraz dla każdego
$\varphi\in V$ i $t\in[0,T]$ spełnia tożsamość
\begin{equation}\label{eq:pde28-weak}
 (u(t),\varphi)_H-(u_0,\varphi)_H
 +\nu\int_0^t(\nabla u(s),\nabla\varphi)_{L^2}\,\mathrm d\lambda_1(s)
 +\int_0^t b(u(s),u(s),\varphi)\,\mathrm d\lambda_1(s)=0.
\end{equation}
Zapis $C([0,T];H_{\mathrm w})$ oznacza ciągłość wszystkich
funkcji skalarnych $t\mapsto(u(t),h)_H$, $h\in H$;
wybieramy ten słabo ciągły reprezentant klasy
funkcji określonych prawie wszędzie.
Wymagamy ponadto nierówności energii
\begin{equation}\label{eq:pde28-energy}
 \frac12\|u(t)\|_2^2+\nu\int_0^t\|\nabla u(s)\|_2^2\,\mathrm d\lambda_1(s)
 \le\frac12\|u_0\|_2^2\qquad (0\le t\le T).
\end{equation}
Wymagamy również silnej nierówności energii:
dla prawie każdego $s\in(0,T)$ i każdego $t\in[s,T]$
\begin{equation}\label{eq:pde28-energy-shift}
 \frac12\|u(t)\|_2^2+
 \nu\int_s^t\|\nabla u(\tau)\|_2^2\,\mathrm d\lambda_1(\tau)
 \le\frac12\|u(s)\|_2^2.
\end{equation}
Zbiór dopuszczalnych chwil $s$ można wybrać
niezależnie od $t$. Nierówność~\eqref{eq:pde28-energy}
odpowiada dodatkowo chwili $s=0$.
\end{definicja}

\begin{lemat}[Oszacowanie nieliniowości w przestrzeni dualnej]
\label{lem:pde28-B-est}
Dla $u,v\in V$ niech $B(u,v)\in V'$ będzie określony przez
$\langle B(u,v),\varphi\rangle=b(u,v,\varphi)$.
Dla $u\in V$ zachodzą oszacowania
\[
 \|B(u,u)\|_{V'}\le
 \begin{cases}
 C\|u\|_2\|\nabla u\|_2,&d=2,\\
 C\|u\|_2^{1/2}\|\nabla u\|_2^{3/2},&d=3.
 \end{cases}
\]
Stała $C$ zależy tylko od wymiaru i wielkości torusa.
\end{lemat}
\begin{proof}
Wykorzystujemy antysymetrię w dwóch ostatnich zmiennych:
$|b(u,u,\varphi)|=|b(u,\varphi,u)|
\le\|u\|_4^2\|\nabla\varphi\|_2$.
W wymiarze $2$ stosujemy nierówność Ladyżenskiej
na torusie:
$\|u\|_4^2\le C\|u\|_2\|u\|_{H^1}$,
a następnie równoważność normy $H^1$ i normy $V$.
W wymiarze $3$ interpolujemy między $L^2$ i $L^6$:
$\|u\|_4\le\|u\|_2^{1/4}\|u\|_6^{3/4}$, ponieważ
$1/4=(1/4)/2+(3/4)/6$.
Z zanurzenia Sobolewa $H^1\hookrightarrow L^6$
i nierówności Poincarégo mamy
$\|u\|_6\le C\|\nabla u\|_2$.
Po podniesieniu nierówności do kwadratu
otrzymujemy obie żądane potęgi.
Po podzieleniu przez $\|\varphi\|_V=\|\nabla\varphi\|_2$
dla $\varphi\ne0$ i wzięciu supremum otrzymujemy tezę.
\end{proof}

\begin{uwaga}[Pochodna czasowa rozwiązania słabego]
Niech $\mathcal A:V\to V'$ będzie operatorem
określonym przez
$\langle\mathcal Au,\varphi\rangle=(u,\varphi)_V$.
Dla każdego rozwiązania z Definicji~\ref{def:pde28-leray}
mamy $\mathcal Au\in L^2(0,T;V')$, a z Lematu
~\ref{lem:pde28-B-est} wynika $B(u,u)\in L^q(0,T;V')$,
gdzie $q=2$ dla $d=2$ i $q=4/3$ dla $d=3$.
Silna mierzalność wynika z ciągłości odwzorowań
$\mathcal A$ i $u\mapsto B(u,u)$ na $V$.
Różniczkując~\eqref{eq:pde28-weak} w sensie
dystrybucyjnym względem czasu dla każdego $\varphi\in V$,
otrzymujemy w $V'$
\[
 u_t=-\nu\mathcal Au-B(u,u)\in L^q(0,T;V').
\]
Własność ta dotyczy wszystkich rozwiązań z definicji,
nie tylko skonstruowanych metodą Galerkina.
\end{uwaga}

\section{Globalne istnienie rozwiązania Leraya--Hopfa}
\begin{twierdzenie}[Istnienie słabego rozwiązania w wymiarach $2$ i $3$]
\label{thm:pde28-existence}
Niech $d\in\{2,3\}$, $\nu>0$, $T>0$ i $u_0\in H$.
Istnieje rozwiązanie Leraya--Hopfa zagadnienia~\eqref{eq:pde28-ns}
na $[0,T]$. Każde takie rozwiązanie spełnia
$\partial_tu\in L^{4/3}(0,T;V')$ dla $d=3$,
a $\partial_tu\in L^2(0,T;V')$ dla $d=2$.
Ponadto $u(t)\to u_0$ silnie w $H$ przy $t\downarrow0$.
\end{twierdzenie}
\begin{proof}
\textbf{Krok 1: układy Galerkina.}
Niech $(w_j)_{j\ge1}$ będzie ortonormalną bazą $H$
złożoną z gładkich, rzeczywistych, bezdywergencyjnych
funkcji własnych operatora Stokesa, z uwzględnieniem krotności.
Oznaczmy $H_m=\operatorname{span}\{w_1,\dots,w_m\}$
i niech $P_m$ będzie rzutem ortogonalnym na $H_m$.
Szukamy $u_m:[0,T]\to H_m$ o współrzędnych
spełniających dla $1\le j\le m$
\begin{equation}\label{eq:pde28-galerkin}
 (u_m',w_j)_H+\nu(\nabla u_m,\nabla w_j)_{L^2}
 +b(u_m,u_m,w_j)=0,\quad u_m(0)=P_m u_0.
\end{equation}
Pisząc $u_m=\sum_{i=1}^m a_iw_i$, dostajemy
układ współrzędnych
\[
 a_j'=-\nu\lambda_j a_j-
 \sum_{i,\ell=1}^m a_i a_\ell b(w_i,w_\ell,w_j),
 \qquad a_j(0)=(u_0,w_j)_H.
\]
Prawa strona jest wielomianem w $(a_1,\ldots,a_m)$,
a więc jest lokalnie lipschitzowska. Twierdzenie
Picarda--Lindelöfa daje jednoznaczne rozwiązanie
na maksymalnym przedziale istnienia.

\textbf{Krok 2: energia i istnienie globalne przybliżeń.}
Mnożymy $j$-te równanie przez współrzędną $u_m$ przy $w_j$
i sumujemy. Wykorzystując $b(u_m,u_m,u_m)=0$, otrzymujemy
\begin{equation}\label{eq:pde28-gal-energy}
 \frac12\|u_m(t)\|_2^2+\nu\int_0^t\|\nabla u_m(s)\|_2^2
 \,\mathrm d\lambda_1(s)=\frac12\|P_m u_0\|_2^2\le\frac12\|u_0\|_2^2.
\end{equation}
Norma $u_m$ nie może eksplodować w skończonym czasie,
a skończony układ ODE można przedłużyć na dowolny $[0,T]$.
Oszacowania są jednostajne względem $m$.

\textbf{Krok 3: pochodne czasowe.}
Rozszerzamy $P_m$ z $H$ na $V'$ przez
$\langle P_mF,\varphi\rangle=\langle F,P_m\varphi\rangle$.
Rozszerzenie jest zgodne z rzutem na $H$, ponieważ
rzut jest samosprzężony względem iloczynu skalarnego $H$.
Rzut $P_m$ jest kontrakcją na $V$ i $V'$:
bazę wybrano z funkcji własnych operatora Stokesa,
które są ortogonalne również dla iloczynu skalarnego
$(\nabla v,\nabla w)_{L^2}$. Wobec tego
$\|P_m\varphi\|_V\le\|\varphi\|_V$;
następnie z definicji normy dualnej
$\|P_m F\|_{V'}=\sup_{\|\varphi\|_V\le1}
|\langle F,P_m\varphi\rangle|\le\|F\|_{V'}$.
Z równania Galerkina oraz lematu~\ref{lem:pde28-B-est} mamy
\[
 \|u_m'\|_{V'}\le\nu\|\nabla u_m\|_2+
 \begin{cases}
 C\|u_m\|_2\|\nabla u_m\|_2,&d=2,\\
 C\|u_m\|_2^{1/2}\|\nabla u_m\|_2^{3/2},&d=3.
 \end{cases}
\]
Prawa strona jest ograniczona odpowiednio w
$L^2(0,T)$ i $L^{4/3}(0,T)$; dla pierwszego składnika
w drugim przypadku korzystamy z $T<\infty$.

\textbf{Krok 4: zwartość w przestrzeni i czasie.}
Z oszacowań i refleksywności odpowiednich
przestrzeni $L^2(0,T;V)$ i $L^q(0,T;V')$
oraz twierdzenia Banacha--Alaoglu (z ośrodkowością
predualu $L^1(0,T;H)$) wydzielamy podciąg taki, że
$u_m\stackrel{*}{\rightharpoonup}u$ w $L^\infty(0,T;H)$,
$u_m\rightharpoonup u$ w $L^2(0,T;V)$ i
$u_m'\rightharpoonup g$ w $L^q(0,T;V')$.
Obie granice $u_m$ są równe przez testowanie
funkcjami czasu i wektorami z $H$.
Testowanie pochodnych przez $\eta(t)\varphi$,
$\eta\in C_c^\infty(0,T)$, $\varphi\in V$,
i całkowanie przez części w czasie daje $g=u_t$
w sensie dystrybucyjnym.

Uzasadnimy zwartość czasoprzestrzenną wprost,
modyfikując argument z Twierdzenia
~\ref{thm:pde20-compact} dla wykładnika $q>1$,
w szczególności $q=4/3$.
Dla $j$ ustalonego współczynniki
$c_{m,j}=(u_m,w_j)_H$ są ograniczone jednostajnie
w czasie i spełniają
\[
 |c_{m,j}(t)-c_{m,j}(s)|
 \le \|u_m'\|_{L^q(0,T;V')}\|w_j\|_V
 |t-s|^{1-1/q}.
\]
Twierdzenie Arzeli--Ascolego i wybór podciągu
diagonalnego zapewniają jednostajną zbieżność
każdego współczynnika na $[0,T]$.
Dla ustalonego $N$ rzuty $P_Nu_m$ zbiegają więc
w $L^2(0,T;H)$. Jednocześnie
\[
 \sup_m\|(I-P_N)u_m\|_{L^2(0,T;H)}^2
 \le\lambda_{N+1}^{-1}
 \sup_m\|u_m\|_{L^2(0,T;V)}^2\longrightarrow0.
\]
Dla dowolnej dokładności wybieramy najpierw $N$,
by oba ogony były małe, a następnie indeksy
podciągu, by różnica rzutów była mała.
Podciąg jest Cauchy'ego w $L^2(0,T;H)$ i jego
granica musi być równa $u$ z poprzednich zbieżności.
Otrzymujemy zatem
\begin{equation}\label{eq:pde28-strong}
 u_m\longrightarrow u\quad\text{silnie w }L^2(0,T;H).
\end{equation}
W szczególności $u_m\otimes u_m\to u\otimes u$ w
$L^1((0,T)\times\T^d)$, gdyż
$\|a\otimes a-b\otimes b\|_1
\le(\|a\|_2+\|b\|_2)\|a-b\|_2$ na czasoprzestrzeni.

\textbf{Krok 5: przejście do granicy w równaniu.}
Niech $\varphi\in\bigcup_{j\ge1}H_j$ i
$\eta\in C_c^\infty(0,T)$. Testujemy~\eqref{eq:pde28-galerkin}
przez $\eta(t)\varphi$ i całkujemy w czasie.
Człon liniowy przechodzi do granicy słabo, a człon nieliniowy
po całkowaniu przez części przestrzennym ma postać
\[
 -\int_0^T\eta(t)\int_{\T^d}
 (u_m\otimes u_m):\nabla\varphi\,\mathrm d\lambda_d(x)
 \,\mathrm d\lambda_1(t).
\]
Ponieważ $\eta\nabla\varphi$ jest ograniczona,
silna zbieżność iloczynów w $L^1$ pozwala przejść
do granicy. Dostajemy
\[
 \int_0^T\bigl(- (u,\varphi)_H\eta'
  +\nu(u,\varphi)_V\eta+b(u,u,\varphi)\eta\bigr)
 \,\mathrm d\lambda_1(t)=0.
\]
Gęstość sumy $H_j$ w $V$ i oszacowanie
$B(u,u)\in L^q(0,T;V')$ wynikające z Lematu
~\ref{lem:pde28-B-est} rozszerzają tę tożsamość
na wszystkie $\varphi\in V$ (każdy ze składników
jest ciągłym funkcjonałem względem $\varphi$
w normie $V$). W kroku następnym uzasadnimy
postać całkową z warunkiem początkowym.

\textbf{Krok 6: warunek początkowy i nierówność energii.}
Dla każdego $\varphi\in V$ funkcje
$t\mapsto(u_m(t),\varphi)_H$ są jednostajnie ograniczone
i jednostajnie Hölderowskie w czasie: z kroku 3,
$|(u_m(t)-u_m(s),\varphi)|\le
C\|\varphi\|_V|t-s|^{1-1/q}$, gdzie $q=4/3$ lub $2$.
Na podciągu z kroku 4 dla każdego $j$
$c_{m,j}\to c_j$ jednostajnie na $[0,T]$.
Z oszacowania energii dla każdego $N$ i $t$
$\sum_{j=1}^N|c_j(t)|^2\le\|u_0\|_2^2$.
Szereg $\sum_j c_j(t)w_j$ określa więc wektor
$U(t)\in H$ o normie co najwyżej $\|u_0\|_2$.
Aproksymując dowolny $h\in H$ przez $P_Nh$,
otrzymujemy jednostajną w czasie zbieżność
$(u_m(t),h)_H\to(U(t),h)_H$: błąd ogona
jest nie większy niż $2\|u_0\|_2\|(I-P_N)h\|_2$.
W szczególności $U\in C([0,T];H_{\mathrm w})$
i $u_m(t)\rightharpoonup U(t)$ w $H$ dla każdego $t$.
Silna zbieżność w $L^2(0,T;H)$ identyfikuje $U=u$
prawie wszędzie; odtąd wybieramy reprezentant $U$.
Ponieważ $u_m(0)=P_m u_0\to u_0$ w $H$,
warunek początkowy jest zachowany.
Z kroku 5 dla każdego $\varphi\in V$
dystrybucyjna pochodna ciągłej funkcji
$t\mapsto(u(t),\varphi)_H$ należy do $L^q(0,T)$
i jest równa $-\nu(u,\varphi)_V-b(u,u,\varphi)$.
Jej reprezentant jest zatem absolutnie ciągły,
a całkowanie od $0$ do dowolnego $t$ daje
~\eqref{eq:pde28-weak}.
Dla ustalonego $t$ stosujemy słabą półciągłość dolną
normy $\|u_m(t)\|_2$ i normy
$\|\nabla u_m\|_{L^2(0,t;L^2)}$ do~\eqref{eq:pde28-gal-energy},
otrzymując~\eqref{eq:pde28-energy} dla każdego $t$.
Ze zbieżności~\eqref{eq:pde28-strong} wydzielamy
dalszy podciąg taki, że $u_m(s)\to u(s)$ w $H$
dla prawie każdego $s\in(0,T)$; na przykład
wybieramy podciąg z sumowalnymi kwadratami
norm błędów w $L^2(0,T;H)$ i stosujemy twierdzenie
Tonellego. Dla każdej takiej chwili $s$ odejmujemy
tożsamość energii Galerkina w chwili $s$ od
tożsamości w dowolnej chwili $t\ge s$.
Prawa strona $\frac12\|u_m(s)\|_2^2$ zbiega
do $\frac12\|u(s)\|_2^2$, natomiast po lewej
stosujemy słabą półciągłość dolną normy $u_m(t)$
oraz normy gradientu ograniczonego do $(s,t)$.
Otrzymujemy~\eqref{eq:pde28-energy-shift}.
Jeden podciąg zapewnia słabą zbieżność
we wszystkich chwilach $t$, więc zbiór
dopuszczalnych $s$ nie zależy od $t$.
Wreszcie $u(t)\rightharpoonup u_0$ w $H$ przy $t\downarrow0$,
a nierówność energii daje
$\limsup_{t\downarrow0}\|u(t)\|_2\le\|u_0\|_2$.
Słaba zbieżność wraz ze zbieżnością norm dowodzi
$u(t)\to u_0$ silnie w $H$.
\end{proof}


\begin{uwaga}[Zwartość paraboliczna a nieliniowość Naviera--Stokesa]
\label{rem:pde-v44-ns-compactness}
Zwartość z kroku czwartego dowodu twierdzenia
\ref{thm:pde28-existence} ma dokładnie tę samą konstrukcję
funkcjonalną co w rozdziale~\ref{chap:pde20}:
$V\hookrightarrow H$ zwarte i $H\hookrightarrow V'$ ciągłe,
$u_m$ ograniczone w $L^2(0,T;V)$ oraz $\partial_tu_m$
ograniczone w $L^q(0,T;V')$, gdzie $q=2$ dla $d=2$
i $q=4/3$ dla $d=3$. Sama słaba zbieżność w przestrzeni
energii nie wystarczałaby do przejścia do granicy
w $u_m\otimes u_m$. Silna zbieżność w $L^2(0,T;H)$
zapewnia bowiem bezpośrednio
\[
 \|u_m\otimes u_m-u\otimes u\|_{L^1((0,T)\times\T^d)}
 \le (\|u_m\|_{L^2}+\|u\|_{L^2})\|u_m-u\|_{L^2}\longrightarrow0.
\]
Nie wnioskujemy stąd o silnej zbieżności gradientów ani
o równości energii w trzech wymiarach.
\end{uwaga}


\begin{uwaga}[Ciśnienie rozwiązania słabego]
Dla gładkiego rozwiązania warunek $\diver u=0$
i równanie~\eqref{eq:pde28-ns} dają
$-\Delta p=\sum_{i,j}\partial_i\partial_j(u_i u_j)$.
Na torusie wybieramy normalizację
$\int_{\T^d}p(t,x)\,\mathrm d\lambda_d(x)=0$.
Dla rozwiązania Leraya--Hopfa wzór ten określa $p$
w sensie dystrybucji przestrzennych i czasowych;
nie należy bez dodatkowych założeń traktować go
jako tożsamości klasycznych funkcji. Dla ustalonego czasu w sensie dystrybucji
przestrzennych rozwiązujemy równanie Poissona dla
$p$ o średniej zero; w bazie Fouriera, dla $k\ne0$,
$\widehat p(k)=-|k|^{-2}\sum_{i,j}k_i k_j
\widehat{(u_i u_j)}(k)$.
Współczynniki iloczynów $u_i u_j$ są dobrze określone,
gdy $u(t)\in L^2$. Ustalamy $\widehat p(0)=0$.
Dla $k\ne0$ możemy zapisać wzór jako
$\widehat p(t,k)=-(2\pi)^{-d}\int_{\T^d}
 |k\cdot u(t,x)|^2|k|^{-2} e^{-ik\cdot x}
 \,\mathrm d\lambda_d(x)$.
W szczególności $|\widehat p(t,k)|\le
(2\pi)^{-d}\|u(t)\|_2^2$ jednostajnie w $k$ i $t$.
Współczynniki są mierzalne względem czasu,
a współczynniki gładkich funkcji testowych
maleją szybciej niż dowolna potęga $|k|$,
jednostajnie na zwartych przedziałach czasu.
Szereg definiuje więc dystrybucję na
$(0,T)\times\T^d$, o zerowym współczynniku dla $k=0$.
Bez normalizacji ciśnienie jest wyznaczone
z dokładnością do dystrybucji zależnej tylko od czasu.

Pozostaje wykazać pełne równanie pędu.
Niech $R=u_t-\nu\Delta u+\diver(u\otimes u)$
w sensie dystrybucji. Dla $k\ne0$ słaba postać
~\eqref{eq:pde28-weak}, testowana sinusami i cosinusami
z amplitudami w $k^\perp$, daje $P_k\widehat R(k)=0$,
gdzie $P_k=I-k\otimes k/|k|^2$.
Z warunku nieściśliwości również
$P_k\widehat u_t(k)=\widehat u_t(k)$ i
$P_k\widehat{\Delta u}(k)=\widehat{\Delta u}(k)$.
Wobec tego
$\widehat R(k)=(I-P_k)\widehat{\diver(u\otimes u)}(k)$.
Wzór na $\widehat p(k)$ daje dokładnie
$ik\widehat p(k)=-\widehat R(k)$.
Dla $k=0$ oba współczynniki znikają na mocy
zerowej średniej $u$ i periodyczności.
Warto zauważyć dodatkową regularność ciśnienia:
z oszacowań $L^4$ w Lemacie~\ref{lem:pde28-B-est}
wynika $u\otimes u\in L^q(0,T;L^2)$
z $q=2$ dla $d=2$ i $q=4/3$ dla $d=3$.
Dla macierzy $G=(u_i u_j)_{i,j}$ mnożnik
$-\sum_{i,j}k_i k_j\widehat G_{ij}(k)/|k|^2$
ma normę co najwyżej $1$ względem normy
Frobeniusa, przez nierówność Cauchy'ego--Schwarza.
Twierdzenie Parsevala daje więc
$\|p(t)\|_2\le\|u(t)\otimes u(t)\|_2=\|u(t)\|_4^2$
dla prawie każdego $t$.
Mnożnik jest ograniczonym operatorem liniowym,
stąd $p\in L^q(0,T;L^2)$ i ma średnią zero.
Zatem $R+\nabla p=0$ w sensie dystrybucji;
otrzymaliśmy pełne równanie~\eqref{eq:pde28-ns}.
\end{uwaga}


\begin{twierdzenie}[Równość energii i ciągłość silna w wymiarze dwa]
\label{thm:pde28-energy-equality-2d}
Niech $\nu>0$, $T>0$, $d=2$ i $u_0\in H$.
Każde rozwiązanie Leraya--Hopfa $u$ na $[0,T]$
należy do $C([0,T];H)$ i spełnia dla wszystkich
$t\in[0,T]$ dokładną równość
\begin{equation}\label{eq:pde28-energy-equality-2d}
 \frac12\|u(t)\|_2^2+
 \nu\int_0^t\|\nabla u(s)\|_2^2\,\mathrm d\lambda_1(s)
 =\frac12\|u_0\|_2^2.
\end{equation}
Ponadto, ponieważ na torusie $2\pi$ funkcje w $H$
mają zerową średnią,
$\|u(t)\|_2^2\le e^{-2\nu t}\|u_0\|_2^2$.
\end{twierdzenie}
\begin{proof}
Z lematu~\ref{lem:pde28-B-est} i warunków
$u\in L^\infty(0,T;H)\cap L^2(0,T;V)$ wynika
$\|B(u,u)\|_{V'}\le C\|u\|_2\|\nabla u\|_2\in
L^2(0,T)$. Niech $\mathcal A:V\to V'$ oznacza rozszerzenie
operatora Stokesa określone przez
$\langle\mathcal Au,\varphi\rangle=
(\nabla u,\nabla\varphi)_{L^2}$.
Wówczas $\|\mathcal Au\|_{V'}\le\|u\|_V$,
a równość słaba daje
$u_t=-\nu\mathcal Au-B(u,u)\in L^2(0,T;V')$.
W trójce Hilberta $V\subset H\subset V'$ mamy
zatem $u\in L^2(0,T;V)$ oraz $u_t\in L^2(0,T;V')$.
Argument z Lematu~\ref{lem:pde20-energy} przenosi się
na obecną trójkę: układ $(w_j)$ jest ortonormalny
w $H$, a $(\lambda_j^{-1/2}w_j)$ w $V$.
W przestrzeni $V'$ mamy, na mocy twierdzenia Riesza,
$\|F\|_{V'}^2=\sum_j\lambda_j^{-1}
|\langle F,w_j\rangle|^2$. Rzuty $P_m$ są więc
kontrakcjami i zbiegają silnie do identyczności
w każdej z trzech przestrzeni.
Zbieżność zmajoryzowana daje $P_mu\to u$ w $L^2(V)$
oraz $P_mu_t\to u_t$ w $L^2(V')$, gdzie pomijamy
w skrótach wspólny przedział czasu $(0,T)$.
Dla $z=P_mu-P_nu$ skończenie wymiarowa reguła
łańcuchowa i wybór chwili o normie nie większej
od średniej kwadratowej dają
\[
 \sup_{0\le t\le T}\|z(t)\|_2^2
 \le T^{-1}\|z\|_{L^2(H)}^2+
 2\|z_t\|_{L^2(V')}\|z\|_{L^2(V)}.
\]
Prawa strona dąży do zera, więc rzuty zbiegają
w $C([0,T];H)$ do reprezentanta $u$.
Jest to ten sam reprezentant co w definicji
rozwiązania Leraya--Hopfa: oba są słabo ciągłe
i równe prawie wszędzie.
Przechodząc do granicy w całkowej regule
łańcuchowej dla $P_mu$ za pomocą nierówności
Höldera w czasie, otrzymujemy absolutną ciągłość
$t\mapsto\|u(t)\|_2^2$ i
$\frac12\frac{\mathrm d}{\mathrm dt}\|u(t)\|_2^2
=\langle u_t(t),u(t)\rangle$ prawie wszędzie.
Równanie zachodzi w $V'$ dla prawie każdego $t$:
korzystamy z tożsamości dystrybucyjnej i ośrodkowości
$V$, aby wybrać wspólny zbiór pełnej miary dla
wszystkich testów. W tych chwilach, w których
również $u(t)\in V$, parujemy równanie z $u(t)$.
Dla takich chwil $b(u,u,u)=0$ na mocy
antysymetrii i gęstości gładkich pól bezdywergencyjnych.
Otrzymujemy pochodną
$\frac12(\|u\|_2^2)'=-\nu\|\nabla u\|_2^2$;
całkowanie w czasie daje~\eqref{eq:pde28-energy-equality-2d}.

Na torusie $\T^2$ najmniejsza dodatnia wartość
własna $-\Delta$ wynosi $1$, więc nierówność
Poincarégo ma postać $\|u\|_2\le\|\nabla u\|_2$.
Dla $E(t)=\|u(t)\|_2^2$ otrzymujemy
$E'(t)\le-2\nu E(t)$ prawie wszędzie.
Całkowanie tej nierówności daje wskazany zanik.
\end{proof}

\begin{wniosek}[Wykładniczy zanik energii rozwiązania Galerkina w 3D]
\label{cor:pde28-decay-3d}
Niech $\nu>0$, $T>0$, $u_0\in H$ i $d=3$.
Rozwiązanie Leraya--Hopfa otrzymane w dowodzie
twierdzenia~\ref{thm:pde28-existence} można wybrać tak, że
\[
 \|u(t)\|_2^2\le e^{-2\nu t}\|u_0\|_2^2,
 \qquad 0\le t\le T.
\]
Nie jest tu zakładana gładkość ani jednoznaczność
rozwiązań trójwymiarowych.
\end{wniosek}
\begin{proof}
Dla przybliżeń Galerkina z~\eqref{eq:pde28-gal-energy}
różniczkujemy energię i używamy
$\|\nabla u_m\|_2^2\ge\|u_m\|_2^2$:
$\frac{\mathrm d}{\mathrm dt}\|u_m(t)\|_2^2
\le-2\nu\|u_m(t)\|_2^2$.
Stąd $\|u_m(t)\|_2^2\le
 e^{-2\nu t}\|P_mu_0\|_2^2\le e^{-2\nu t}\|u_0\|_2^2$.
W kroku 6 dowodu~\ref{thm:pde28-existence} wybrano
jeden podciąg zbieżny słabo w $H$ dla każdej chwili
$t\in[0,T]$. Słaba półciągłość dolna normy przy
ustalonym $t$ daje pożądane oszacowanie dla tej granicy.
W szczególności w tym wniosku nie korzystamy z
nieudowodnionej równości energii dla wszystkich
rozwiązań Leraya--Hopfa w wymiarze trzy.
\end{proof}


\section{Wymiar dwa: jednoznaczność i wirowość}
\begin{twierdzenie}[Jednoznaczność słabych rozwiązań w wymiarze dwa]
\label{thm:pde28-unique2d}
Niech $\nu>0$, $T>0$ i $u_0\in H$ dla $d=2$.
Rozwiązanie Leraya--Hopfa na $[0,T]$ jest jednoznaczne.
Jeśli $u,v$ odpowiadają danym $u_0,v_0\in H$, to
\[
 \|u(t)-v(t)\|_2^2\le\|u_0-v_0\|_2^2
 \exp\left(\frac C\nu\int_0^t\|\nabla v(s)\|_2^2\,\mathrm d\lambda_1(s)\right),
 \qquad 0\le t\le T.
\]
\end{twierdzenie}
\begin{proof}
Z oszacowania w Lemacie~\ref{lem:pde28-B-est}
i słabej postaci równań wynika
$u_t,v_t\in L^2(0,T;V')$ dla wszystkich rozwiązań
Leraya--Hopfa w wymiarze dwa.
Stosujemy regułę łańcuchową w trójce Hilberta,
uzasadnioną w dowodzie Twierdzenia
~\ref{thm:pde28-energy-equality-2d}, do $w=u-v$.
Odejmujemy równania w $V'$ i parujemy z $w(t)$
dla prawie każdego $t$.
Z trójliniowości mamy
$b(u,u,w)-b(v,v,w)=b(u,w,w)+b(w,v,w)$;
pierwszy składnik znika przez antysymetrię.
Otrzymujemy
\[
 \frac12\frac{\mathrm d}{\mathrm dt}\|w\|_2^2
 +\nu\|\nabla w\|_2^2=-b(w,v,w).
\]
Wymiar dwa pozwala użyć nierówności Ladyżenskiej:
$|b(w,v,w)|\le\|w\|_4^2\|\nabla v\|_2
\le C\|w\|_2\|\nabla w\|_2\|\nabla v\|_2$.
Nierówność Younga $ab\le\nu a^2/2+b^2/(2\nu)$
dla $a=\|\nabla w\|_2$,
$b=C\|w\|_2\|\nabla v\|_2$ daje
\[
 \frac{\mathrm d}{\mathrm dt}\|w\|_2^2
 +\nu\|\nabla w\|_2^2
 \le (C^2/\nu)\|\nabla v\|_2^2\|w\|_2^2.
\]
Po pominięciu składnika nieujemnego i zmianie
stałej otrzymujemy żądaną nierówność różniczkową.
Ponieważ $\nabla v\in L^2(0,T;L^2)$,
lemat Gronwalla daje podane oszacowanie, a dla równych
danych otrzymujemy $u=v$.
\end{proof}

\begin{stwierdzenie}[Równanie wirowości i zanik enstrofii]
\label{prop:pde28-vorticity}
Niech $d=2$, a $(u,p)$ będzie gładkim rozwiązaniem
\eqref{eq:pde28-ns} na $[0,T]\times\T^2$.
Przyjmijmy $\omega=\partial_1u_2-\partial_2u_1$
oraz $\mathcal E_\omega(t)=\frac12\|\omega(t)\|_2^2$.
Wówczas
\[
 \partial_t\omega+u\cdot\nabla\omega=\nu\Delta\omega,
 \qquad \mathcal E_\omega(t)+\nu\int_0^t\|\nabla\omega(s)\|_2^2
 \,\mathrm d\lambda_1(s)=\mathcal E_\omega(0).
\]
Dodatkowo dla bezdywergencyjnych pól o zerowej średniej
$\|\nabla u(t)\|_2=\|\omega(t)\|_2$.
\end{stwierdzenie}
\begin{proof}
Stosując operator $\partial_1$ do drugiego równania
wektorowego i odejmując $\partial_2$ pierwszego,
eliminujemy gradient ciśnienia, ponieważ pochodne
mieszane są równe. Dla członu konwekcyjnego
stosujemy regułę Leibniza:
\[
 \begin{aligned}
 &\partial_1((u\cdot\nabla)u_2)
 -\partial_2((u\cdot\nabla)u_1)\\
 &\quad=u\cdot\nabla\omega+
 (\partial_1u_1+\partial_2u_2)
 (\partial_1u_2-\partial_2u_1).
 \end{aligned}
\]
Ostatni iloczyn znika na mocy nieściśliwości.
Operator rotacji komutuje z laplasjanem,
co daje równanie wirowości.
Mnożymy równanie przez $\omega$ i całkujemy po torusie.
Składnik konwekcyjny znika:
$\int_{\T^2}u\cdot\nabla(\omega^2/2)\,\mathrm d\lambda_2(x)=0$.
Człon dyfuzyjny daje $-\nu\|\nabla\omega\|_2^2$
po periodycznym całkowaniu przez części.
Całkowanie tożsamości różniczkowej w czasie daje
równość enstrofii.
Równość norm można wykazać na współczynnikach
Fouriera: dla $k\ne0$
\[
 |k|^2|\widehat u(k)|^2
 =|k\cdot\widehat u(k)|^2+
 |k_1\widehat u_2(k)-k_2\widehat u_1(k)|^2.
\]
Pierwszy składnik po prawej znika, a drugi to
$|\widehat\omega(k)|^2$.
Sumowanie i twierdzenie Parsevala dają
$\|\nabla u\|_2=\|\omega\|_2$.
Po pomnożeniu równości przez $|k|^2$ otrzymujemy
także $\|\Delta u\|_2=\|\nabla\omega\|_2$
dla pól w $H^2$.
\end{proof}

\begin{twierdzenie}[Globalne oszacowanie $H^1$ w wymiarze dwa]
\label{thm:pde28-H1-2d}
Niech $d=2$, $\nu>0$, $T>0$ oraz $u_0\in V$.
Jedyne rozwiązanie Leraya--Hopfa $u$ na $[0,T]$ spełnia
$u\in C([0,T];V)\cap L^2(0,T;D(A))$
i $u_t\in L^2(0,T;H)$
oraz dla każdego $t\in[0,T]$
\[
 \|\nabla u(t)\|_2^2+2\nu\int_0^t\|\Delta u(s)\|_2^2
 \,\mathrm d\lambda_1(s)\le\|\nabla u_0\|_2^2.
\]
\end{twierdzenie}
\begin{proof}
Wybieramy w dowodzie twierdzenia~\ref{thm:pde28-existence}
przestrzenie Galerkina złożone ze wszystkich bezdywergencyjnych
modów Fouriera $0<|k|\le N$ (są to podprzestrzenie
odpowiadające pełnym grupom wartości własnych).
Oznaczmy przez $Q_N$ obcięcie do tych częstości,
działające składowo na polach wektorowych.
Rzut Galerkina jest równy $Q_NP$ dla pól o średniej zero.
Równanie przybliżone ma postać
$\partial_t u_N-\nu\Delta u_N+
Q_NP((u_N\cdot\nabla)u_N)=0$.
Rotacja komutuje z $Q_N$ i $\Delta$, a usuwa
część gradientową $I-P$ (na każdej częstości
jest ona równoległa do $k$).
Dla $\omega_N=\partial_1u_{N,2}
-\partial_2u_{N,1}$ z obliczenia w Stwierdzeniu
~\ref{prop:pde28-vorticity} równanie rzutowane ma postać
$\partial_t\omega_N+Q_N(u_N\cdot\nabla\omega_N)
=\nu\Delta\omega_N$, gdzie $Q_N$ jest ortogonalnym
rzutem na skalarne mody z $0<|k|\le N$.
Ponieważ $Q_N\omega_N=\omega_N$, mnożenie przez
$\omega_N$ daje
\[
 \frac12\frac{\mathrm d}{\mathrm dt}\|\omega_N\|_2^2
 +\nu\|\nabla\omega_N\|_2^2=0.
\]
Przy tym człon transportowy znika przez $\diver u_N=0$.
Dla bezdywergencyjnego pola o zerowej średniej zachodzą
$\|\omega_N\|_2=\|\nabla u_N\|_2$ i
$\|\nabla\omega_N\|_2=\|\Delta u_N\|_2$,
co można sprawdzić oddzielnie na każdym modzie Fouriera.
Otrzymujemy zatem jednostajne ograniczenia w
$L^\infty(0,T;V)$ i $L^2(0,T;H^2)$.
Dokładniej, $\|\nabla u_N(0)\|_2\le\|\nabla u_0\|_2$,
a $\|\Delta u_N\|_2$ kontroluje normę $H^2$
dzięki zerowej średniej i oszacowaniu
$(1+|k|^2)^2\le4|k|^4$ dla $k\ne0$.
Wydzielamy podciąg słabo zbieżny w $L^2(H^2)$
i słabo-* zbieżny w $L^\infty(V)$.
Obie granice są równe granicy z dowodu istnienia
dzięki silnej zbieżności w $L^2(H)$.
Dla każdej chwili $t$ zachodzi
$u_N(t)\rightharpoonup u(t)$ w $V$:
istotnie, ciąg jest ograniczony w $V$,
a jego słaba granica w $H$ została już ustalona
dla każdej chwili; każda słabo zbieżna
w $V$ podsekwencja ma zatem tę samą granicę.
W szczególności $u(t)\in V$ dla każdego $t$.
Stosując słabą półciągłość dolną normy $V$
oraz normy $\Delta u_N$ w $L^2(0,t;L^2)$
do tożsamości enstrofii, otrzymujemy żądane
oszacowanie nawet dla wszystkich $t\in[0,T]$.
Na mocy jednoznaczności z Twierdzenia
~\ref{thm:pde28-unique2d} jest to jedyne
rozwiązanie Leraya--Hopfa.

Możemy również uzasadnić silną ciągłość w $V$.
Dla prawie każdego $t$ pole $u(t)\in D(A)$,
a nierówność Ladyżenskiej zastosowana do
pochodnych pierwszego rzędu (które mają średnią zero)
daje
\[
 \|\nabla u\|_4^2\le C\|\nabla u\|_2\|\Delta u\|_2.
\]
Łącząc to z $\|u\|_4^2\le C\|u\|_2\|\nabla u\|_2$,
otrzymujemy
\[
 \|P((u\cdot\nabla)u)\|_2^2
 \le C\|u\|_2\|\nabla u\|_2^2\|\Delta u\|_2.
\]
Prawa strona jest całkowalna po czasie,
gdyż $u$ jest ograniczona w $V$ i
$\Delta u\in L^2(0,T;H)$. Z równania rzutowanego
wynika więc $u_t\in L^2(0,T;H)$.
Niech $A^{1/2}$ będzie mnożnikiem
$\lambda_j^{1/2}$ w bazie Stokesa.
Dla $z=A^{1/2}u$ mamy $z\in L^2(0,T;V)$,
gdyż $\|z\|_V^2=\|Au\|_2^2$, oraz
$z_t=A^{1/2}u_t\in L^2(0,T;V')$, ponieważ
$\|A^{1/2}h\|_{V'}=\|h\|_2$ dla $h\in H$.
Równość pochodnych sprawdzamy na skończonych
rzutach spektralnych, a następnie przechodzimy
do granicy w obu przestrzeniach Bochnera.
Reguła łańcuchowa z dowodu Twierdzenia
~\ref{thm:pde28-energy-equality-2d} daje
$z\in C([0,T];H)$. Ponieważ $A^{-1/2}:H\to V$
jest izometrią, otrzymujemy $u\in C([0,T];V)$;
reprezentant ten jest zgodny z wcześniejszym
ciągłym reprezentantem w $H$.
\end{proof}

\begin{uwaga}[Dlaczego wymiar trzy jest inny]
W wymiarze $3$ wirowość $\boldsymbol\omega=\nabla\times u$
spełnia dla gładkiego rozwiązania
$\partial_t\boldsymbol\omega+(u\cdot\nabla)\boldsymbol\omega
=(\boldsymbol\omega\cdot\nabla)u+\nu\Delta\boldsymbol\omega$.
Pojawia się człon \emph{rozciągania wirów}, który nie ma
odpowiednika w wymiarze $2$. Z tożsamości energii $L^2$
nie otrzymuje się wprost ograniczenia normy $\|\nabla u\|_2$
we wszystkich czasach w wymiarze $3$.
\end{uwaga}

\section{Wymiar trzy: kryterium warunkowej jednoznaczności}
\begin{twierdzenie}[Jednoznaczność przy dodatkowej całkowalności]
\label{thm:pde28-serrin}
Niech $d=3$, $\nu>0$, $T>0$ i $u_0\in H$.
Niech $u,v$ będą rozwiązaniami Leraya--Hopfa na $[0,T]$
o tych samych danych początkowych. Załóżmy dodatkowo, że
$v\in L^\infty(0,T;V)\cap L^2(0,T;D(A))$,
$\partial_tv\in L^2(0,T;H)$ oraz
$v\in L^q(0,T;L^p(\T^3))$, gdzie $p>3$ i
$2/q+3/p\le1$. Wówczas $u=v$ prawie wszędzie na
$(0,T)\times\T^3$.
\end{twierdzenie}
\begin{proof}
Z Lematu~\ref{lem:pde28-B-est} i słabej postaci
równania wiemy, że $u_t\in L^{4/3}(0,T;V')$.
Ponieważ $v\in L^\infty(0,T;V)$ i
$v_t\in L^2(0,T;H)$, reguła łańcuchowa uzasadniona
w dowodzie Twierdzenia~\ref{thm:pde28-energy-equality-2d}
daje $v\in C([0,T];H)$ oraz równość energii dla $v$:
parujemy równanie dla $v$ z $v$ w $V'$ i stosujemy
$b(v,v,v)=0$. Nie stosujemy tej reguły bezpośrednio
do $u$, gdyż jego pochodna czasowa nie musi należeć
do $L^2(0,T;V')$.

Aby uzyskać tożsamość mieszaną, niech $P_n$ oznacza
rzut na pierwsze $n$ wektorów bazy Stokesa.
Współczynniki $P_nu$ i $P_nv$ są absolutnie ciągłe:
ich pochodne są odpowiednio współczynnikami $u_t$
i $v_t$. Z reguły Leibniza w skończonym wymiarze
otrzymujemy dla każdego $t\in[0,T]$
\[
 \begin{aligned}
 &(P_nu(t),P_nv(t))_H-(P_nu_0,P_nv(0))_H\\
 &\quad=\int_0^t
 \bigl(\langle u_t,P_nv\rangle+(P_nu,v_t)_H\bigr)
 \,\mathrm d\lambda_1(s).
 \end{aligned}
\]
Po lewej stronie przechodzimy do granicy w $H$
dla każdej ustalonej chwili, korzystając z
wybranych słabo ciągłych reprezentantów.
Po prawej stronie $P_nv\to v$ w $L^4(0,T;V)$
oraz $P_nu\to u$ w $L^2(0,T;H)$ na mocy
silnej zbieżności rzutów i zbieżności zmajoryzowanej.
Nierówności Höldera dają zbieżność obu całek,
jednostajną względem górnej granicy całkowania.
Stąd
\[
 (u(t),v(t))_H-(u_0,v(0))_H
 =\int_0^t
 (\langle u_t,v\rangle+(u,v_t)_H)\,\mathrm d\lambda_1(s).
\]
Oba równania zachodzą w $V'$ prawie wszędzie.
Wstawiając je do ostatniej tożsamości, dostajemy
\[
 \begin{aligned}
 (u(t),v(t))_H={}&(u_0,v(0))_H\\
 &-2\nu\int_0^t(u,v)_V\,\mathrm d\lambda_1(s)\\
 &-\int_0^t
 [b(u,u,v)+b(v,v,u)]\,\mathrm d\lambda_1(s).
 \end{aligned}
\]
Całkowalność pierwszego członu nieliniowego
wynika z $B(u,u)\in L^{4/3}(V')$ i $v\in L^4(V)$;
drugiego z $B(v,v)\in L^\infty(V')$
oraz $u\in L^2(V)$. W tych skrótach przedziałem
całkowania jest $(0,T)$.
Łącząc otrzymaną tożsamość mieszaną z nierównością
energii dla $u$ i równością energii dla $v$,
uzyskujemy dla $w=u-v$ nierówność
\[
 \frac12\|w(t)\|_2^2+\nu\int_0^t\|\nabla w\|_2^2\,\mathrm d\lambda_1(s)
 \le-\int_0^t b(w(s),v(s),w(s))\,\mathrm d\lambda_1(s).
\]
Istotnie, sumujemy nierówność energii dla $u$
i równość energii dla $v$, po czym odejmujemy
tożsamość mieszaną. Składniki lepkościowe łączą się
w $\nu\|\nabla(u-v)\|_2^2$, a składnik początkowy
jest równy $\frac12\|u_0-v(0)\|_2^2=0$.
Po prawej stronie pozostaje
$b(u,u,v)+b(v,v,u)=-b(w,v,w)$:
wstawiamy $u=v+w$ i stosujemy antysymetrię,
w szczególności $b(a,v,v)=0$ oraz
$b(w,w,v)=-b(w,v,w)$.
Otrzymujemy zatem nierówność względnej energii.

Niech najpierw $3<p<\infty$.
Z nierówności Höldera i interpolacji pomiędzy
$L^2$ i $L^6$, przy parametrze $\theta=3/p$
spełniającym $1/(2p/(p-2))=(1-\theta)/2+\theta/6$,
mamy
\[
 |b(w,w,v)|\le\|v\|_p\|w\|_{2p/(p-2)}\|\nabla w\|_2
 \le C\|v\|_p\|w\|_2^{1-3/p}
 \|\nabla w\|_2^{1+3/p}.
\]
Stosujemy nierówność Younga z wykładnikami
$2/(1+\theta)$ i $2/(1-\theta)$,
dobierając współczynnik przy pierwszym składniku
tak, aby pochłonąć połowę dyssypacji.
Dla $r=2/(1-\theta)=2p/(p-3)$ dostajemy
\[
 |b(w,w,v)|\le\frac{\nu}{2}\|\nabla w\|_2^2+
 C_{\nu,p}\|v\|_p^r\|w\|_2^2.
\]
Dla $p=\infty$ to samo oszacowanie obowiązuje
z $r=2$: bez interpolacji stosujemy
$|b(w,w,v)|\le\|v\|_\infty\|w\|_2\|\nabla w\|_2$
i nierówność Younga.
Warunek $2/q+3/p\le1$ zapewnia $q\ge r$,
więc $\|v\|_p^r$ jest całkowalna w czasie.
Z antysymetrii i nierówności względnej energii
wynika zatem
\[
 \|w(t)\|_2^2\le
 2C_{\nu,p}\int_0^t\|v(s)\|_p^r\|w(s)\|_2^2
 \,\mathrm d\lambda_1(s).
\]
Funkcja po lewej stronie jest mierzalna
i ograniczona. Całkowa postać lematu Gronwalla
daje $\|w(t)\|_2=0$ dla wszystkich $t\in[0,T]$.
\end{proof}

\begin{uwaga}[Co pozostaje otwarte]
Dla $d=3$ twierdzenie~\ref{thm:pde28-existence} gwarantuje
istnienie globalnych rozwiązań słabych, ale nie dowodzi
ich jednoznaczności ani gładkości dla dowolnych danych w $H$.
Dowód twierdzenia~\ref{thm:pde28-serrin} jest dowodem
jednoznaczności słabego rozwiązania w obecności rozwiązania
o silniejszej regularności. Przy podanych założeniach
warunek całkowalności Prodiego--Serrina jest w istocie
automatyczny dla pewnej pary wykładników:
$v\in L^\infty(0,T;V)$ i zanurzenie Sobolewa dają
$v\in L^\infty(0,T;L^6)$, a przy $q=\infty$ mamy $2/q+3/6=1/2<1$.
Zatem do jednoznaczności wystarczają tu same
podane założenia o silniejszej regularności.
Zachowaliśmy rachunek dla ogólnych wykładników
$p,q$, by wyjaśnić znaczenie warunku
$2/q+3/p\le1$. Usunięcie tych dodatkowych założeń
i pozostawienie jedynie warunku $v\in L^q_tL^p_x$
daje znacznie silniejsze twierdzenie Prodiego--Serrina,
którego pełnego dowodu tutaj nie podajemy.
Argument wymiaru dwa wykorzystuje całkowalność
$\|\nabla v\|_2^2$ i nierówność Ladyżenskiej, której
odpowiednik w wymiarze trzy daje zbyt wysoką potęgę
$\|\nabla v\|_2$. Przejście od warunkowej jednoznaczności
do pełnej wymagałoby wykazania dodatkowej regularności;
nie należy tego zakładać bez dowodu.
\end{uwaga}

\section{Lokalne oszacowanie silnych rozwiązań i skala problemu}
\begin{stwierdzenie}[Oszacowanie normy $H^1$ w wymiarze trzy]
\label{prop:pde28-H1}
Niech $d=3$, $\nu>0$ i $(u,p)$ będzie gładkim rozwiązaniem
\eqref{eq:pde28-ns} na $[0,T]\times\T^3$, o średniej $u$ równej zero.
Istnieje $C>0$, zależna tylko od torusa, taka że
\[
 \frac{\mathrm d}{\mathrm dt}\|\nabla u(t)\|_2^2
 +\nu\|\Delta u(t)\|_2^2
 \le C\nu^{-3}\|\nabla u(t)\|_2^6.
\]
\end{stwierdzenie}
\begin{proof}
Bierzemy iloczyn skalarny równania z $-\Delta u$.
Ponieważ $\diver\Delta u=\Delta\diver u=0$,
człon ciśnienia znika przez całkowanie przez części.
Człon czasowy daje
$(u_t,-\Delta u)_H=(\nabla u_t,\nabla u)_{L^2}
=\frac12(\|\nabla u\|_2^2)'$, a lepkościowy
$\nu\|\Delta u\|_2^2$. Stosując nierówność Höldera
z wykładnikami $6,3,2$ do członu konwekcyjnego,
otrzymujemy
\[
 \frac12\frac{\mathrm d}{\mathrm dt}\|\nabla u\|_2^2
 +\nu\|\Delta u\|_2^2
 \le\|u\|_6\|\nabla u\|_3\|\Delta u\|_2.
\]
Z zanurzenia Sobolewa i zerowej średniej mamy
$\|u\|_6\le C\|\nabla u\|_2$.
Każda pochodna przestrzenna $u$ również ma
średnią zero, więc $\|\nabla u\|_6\le C\|D^2u\|_2$.
W bazie Fouriera suma kwadratów wszystkich
pochodnych drugiego rzędu (z uporządkowanymi
indeksami) ma mnożnik
$\sum_{i,j}k_i^2 k_j^2=|k|^4$, stąd
$\|D^2u\|_2=\|\Delta u\|_2$.
Interpolacja $\|\nabla u\|_3\le
\|\nabla u\|_2^{1/2}\|\nabla u\|_6^{1/2}$
daje więc
$\|\nabla u\|_3\le C\|\nabla u\|_2^{1/2}
\|\Delta u\|_2^{1/2}$.
Stąd prawa strona jest ograniczona przez
$C\|\nabla u\|_2^{3/2}\|\Delta u\|_2^{3/2}$.
Nierówność Younga z wykładnikami $4/3$ i $4$
pozwala pochłonąć połowę składnika
$\nu\|\Delta u\|_2^2$ i pozostawia
$C\nu^{-3}\|\nabla u\|_2^6$; pomnożenie przez $2$
i zmianę stałej kończy dowód.
\end{proof}

\begin{uwaga}[Czas obowiązywania oszacowania a priori]
Jeśli $y(t)=\|\nabla u(t)\|_2^2$ oraz $y_0=y(0)>0$,
to Stwierdzenie~\ref{prop:pde28-H1} daje
$y'\le C\nu^{-3}y^3$. Z porównania ze skalarnym
równaniem $z'=C\nu^{-3}z^3$, $z(0)=y_0$,
wynika
\[
 y(t)\le
 \frac{y_0}{\sqrt{1-2C\nu^{-3}y_0^2t}},
 \qquad 0\le t<\frac{\nu^3}{2Cy_0^2},
\]
o ile gładkie rozwiązanie istnieje w rozważanej chwili.
Wzór na $z$ wynika z różniczkowania $z^{-2}$;
porównanie stosujemy na każdym zwartym przedziale
wspólnego istnienia, gdzie sześcian jest lipschitzowski
na wspólnym ograniczonym zakresie.
Dla $y_0=0$ zerowa średnia daje $u_0=0$,
a tożsamość energii wymusza $u=0$.
Powyższy rachunek jest oszacowaniem a priori,
a nie samodzielnym dowodem istnienia silnego
rozwiązania ani dowodem eksplozji w chwili,
w której mianownik oszacowania znika.
\end{uwaga}

\begin{kontrprzyklad}[Dlaczego sama energia $L^2$ nie wystarcza]
Funkcje $u_n(x)=\sin(nx_1)e_2$ na $\T^3$, gdzie
$e_2=(0,1,0)$, mają zerową dywergencję i średnią.
Ich normy $\|u_n\|_2$ są jednakowo ograniczone,
podczas gdy $\|\nabla u_n\|_2\to\infty$.
Wobec tego kontrola samej energii $L^2$ nie daje
jednostajnego oszacowania $H^1$ w pojedynczych chwilach.
Nie jest to kontrprzykład dla regularności rozwiązań
równań Naviera--Stokesa, lecz przykład braku
odpowiedniego wniosku funkcjonalno-analitycznego.
\end{kontrprzyklad}

\begin{uwaga}[Skalowanie na całej przestrzeni i na torusie]
Dla $\ell>0$ połóżmy
\[
 u_\ell(t,x)=\ell u(\ell^2t,\ell x),\qquad
 p_\ell(t,x)=\ell^2p(\ell^2t,\ell x).
\]
Piszemy $s=\ell^2t$, $y=\ell x$. Z reguły
łańcuchowej otrzymujemy
$\partial_t u_\ell=\ell^3\partial_s u$,
$\nabla_xu_\ell=\ell^2\nabla_yu$,
$\Delta_xu_\ell=\ell^3\Delta_yu$ oraz
$\nabla_xp_\ell=\ell^3\nabla_yp$,
gdzie prawe strony obliczamy w $(s,y)$.
W szczególności $(u_\ell\cdot\nabla_x)u_\ell
=\ell^3(u\cdot\nabla_y)u$,
a $\diver_x u_\ell=\ell^2\diver_y u$.
Dlatego przeskalowane pole spełnia równanie
Naviera--Stokesa z tą samą lepkością $\nu$.
Na $\R^d$ zmiana zmiennych daje
\[
 \|u_\ell(t)\|_{L^p(\R^d)}
 =\ell^{1-d/p}\|u(\ell^2t)\|_{L^p(\R^d)}
 \quad(1\le p\le\infty),
\]
gdzie $d/\infty=0$. Dla normy mieszanej na
odpowiednio przeskalowanym przedziale czasu
czynnik wynosi $\ell^{1-d/p-2/q}$, również
dla $q=\infty$ z konwencją $2/\infty=0$.
W wymiarze trzy warunek $2/q+3/p=1$
jest więc warunkiem niezmienniczości tej normy.
Norma $L^2$ skaluje się z czynnikiem
$\ell^{1-d/2}$: jest niezmiennicza w wymiarze dwa,
ale nie w wymiarze trzy.
Na torusie o okresie $2\pi$ powyższa transformacja
prowadzi naturalnie do torusa o okresie $2\pi/\ell$;
nie jest dowolnym ciągłym skalowaniem
zagadnienia na ustalonym torusie.
\end{uwaga}

\section{Dokładne rozwiązanie: wir Taylora--Greena}
\begin{przyklad}[Wir Taylora--Greena na torusie dwuwymiarowym]
\label{ex:pde28-TG}
Niech $\nu>0$ i $(x,y)\in\T^2$. Funkcje
\begin{align*}
 u(t,x,y)&=e^{-2\nu t}(\sin x\cos y,-\cos x\sin y),\\
 p(t,x,y)&=\tfrac14e^{-4\nu t}(\cos 2x+\cos 2y)
\end{align*}
rozwiązują~\eqref{eq:pde28-ns} z danymi
$u_0=(\sin x\cos y,-\cos x\sin y)$.
Oznaczmy $a=e^{-2\nu t}$. Bezpośrednio
$\partial_xu_1=a\cos x\cos y=-\partial_yu_2$,
$\Delta u=-2u$ i $u_t=-2\nu u$.
Rozpisując konwekcję, otrzymujemy
\[
 \begin{aligned}
 ((u\cdot\nabla)u)_1
 &=a^2\sin x\cos x(\cos^2y+\sin^2y),\\
 ((u\cdot\nabla)u)_2
 &=a^2\sin y\cos y(\sin^2x+\cos^2x).
 \end{aligned}
\]
Jest to $\frac12a^2(\sin2x,\sin2y)=-\nabla p$.
Średnie $u$ i $p$ są zerowe. Ponieważ
$\int_0^{2\pi}\sin^2x\,\mathrm d\lambda_1(x)
=\int_0^{2\pi}\cos^2x\,\mathrm d\lambda_1(x)=\pi$,
mamy $\|u\|_2^2=2\pi^2a^2$.
Z całkowania przez części
$\|\nabla u\|_2^2=(u,-\Delta u)_H=2\|u\|_2^2$,
a różniczkowanie daje $\omega=2a\sin x\sin y$.
Ponadto
\[
 \frac12\|u(t)\|_2^2=\pi^2e^{-4\nu t},\qquad
 \|\nabla u(t)\|_2^2=4\pi^2e^{-4\nu t},\qquad
 \omega(t,x,y)=2e^{-2\nu t}\sin x\sin y.
\]
Różniczkowanie energii daje
$-4\nu\pi^2e^{-4\nu t}=-\nu\|\nabla u(t)\|_2^2$.
Jawny wzór pozwala porównać obliczenia numeryczne
z rozwiązaniem dokładnym na dowolnej siatce
i badać błąd oraz rząd zbieżności schematu.

\begin{center}
\begin{tikzpicture}[x=1.5cm,y=1.5cm,>=Stealth,font=\footnotesize]
 \foreach \i in {0,...,15}{
  \foreach \j in {0,...,15}{
   \pgfmathsetmacro{\xx}{360*(\i+.5)/16}
   \pgfmathsetmacro{\yy}{360*(\j+.5)/16}
   \pgfmathsetmacro{\ww}{sin(\xx)*sin(\yy)}
   \pgfmathsetmacro{\opa}{.65*abs(\ww)}
   \ifdim\ww pt>0pt
    \fill[red,opacity=\opa] ({\i/4},{\j/4})
       rectangle ({(\i+1)/4},{(\j+1)/4});
   \else
    \fill[blue,opacity=\opa] ({\i/4},{\j/4})
       rectangle ({(\i+1)/4},{(\j+1)/4});
   \fi
  }
 }
 \foreach \i in {0,...,7}{
  \foreach \j in {0,...,7}{
   \pgfmathsetmacro{\xx}{360*(\i+.5)/8}
   \pgfmathsetmacro{\yy}{360*(\j+.5)/8}
   \draw[-{Stealth[length=2.5pt]},black!85]
   ({(\i+.5)/2},{(\j+.5)/2})
   --++({.23*sin(\xx)*cos(\yy)},{-.23*cos(\xx)*sin(\yy)});
  }
 }
 \draw[->] (0,0)--(4.3,0) node[right] {$x$};
 \draw[->] (0,0)--(0,4.3) node[above] {$y$};
 \draw[black!60] (0,0) rectangle (4,4);
 \foreach \p/\lab in {0/0,2/\pi,4/{2\pi}}{
  \node[below] at (\p,0) {$\lab$};
  \node[left] at (0,\p) {$\lab$};
 }
\end{tikzpicture}

{\small Rysunek: pole prędkości i wirowość w chwili $t=0$.
Czerwony kolor oznacza dodatnią, niebieski ujemną wirowość;
strzałki są obliczone z jawnego wzoru.}
\end{center}
\end{przyklad}


\begin{przyklad}[Trójwymiarowy wir Taylora--Greena: dane początkowe]
\label{ex:pde28-TG3D}
Na torusie $\T^3=([0,2\pi]/\sim)^3$ rozważmy
bezdywergencyjne pole początkowe
\[
 u_0(x,y,z)=(\sin x\cos y\cos z,\,-\cos x\sin y\cos z,\,0).
\]
Istotnie, $\partial_x u_{0,1}+\partial_y u_{0,2}=0$,
a $\partial_z u_{0,3}=0$. Wszystkie składowe
mają średnią zero, więc $u_0\in H\cap C^\infty$.
Jego wirowość jest \emph{wektorem}
\[
 \nabla\times u_0=(-\cos x\sin y\sin z,\,
 -\sin x\cos y\sin z,\,2\sin x\sin y\cos z).
\]
W szczególności trzecia składowa wirowości wynosi
$\omega_{0,z}=2\sin x\sin y\cos z$. Energia początkowa
względem miary Lebesgue'a na sześcianie okresowym wynosi
\[
 E(0)=\frac12\int_{[0,2\pi]^3}|u_0(x,y,z)|^2
 \,\mathrm d\lambda_3(x,y,z)=\pi^3.
\]
Istotnie, całka z kwadratu każdej z dwóch
niezerowych składowych jest równa $\pi^3$.
Ponadto $\Delta u_0=-3u_0$ i
$\|\nabla u_0\|_2^2=6\pi^3$.
Dla gładkiej ewolucji z tymi danymi
początkowa pochodna energii wynosi więc $-6\nu\pi^3$.

Ponieważ $u_{0,3}=0$, \emph{linie prądu pola początkowego}
leżą w płaszczyznach $z=\mathrm{const}$. Dla ustalonego
$z$ regularne linie poziomicowe funkcji prądu
$\psi_z(x,y)=\sin x\sin y\cos z$ są liniami prądu,
ponieważ $(u_{0,1},u_{0,2})=(\partial_y\psi_z,
-\partial_x\psi_z)$, skąd $u_0\cdot\nabla\psi_z=0$
przy ustalonym $z$. W punktach stagnacji
trajektorie są stałe; dla $\cos z=0$ całe pole zanika.

Dane początkowe nie są jawnym rozwiązaniem
trójwymiarowego zagadnienia w czasie.
W istocie bezpośrednio
\[
 (u_0\cdot\nabla)u_0
 =\tfrac12\cos^2z(\sin2x,\sin2y,0).
\]
Równanie Poissona na ciśnienie w chwili początkowej
ma prawą stronę
$(\cos2x+\cos2y)\cos^2z$.
Jedynym rozwiązaniem o średniej zero jest
\[
 p_0=\tfrac1{16}(\cos2x+\cos2y)(2+\cos2z),
\]
co sprawdzamy, stosując $-\Delta$ do obu składników
i używając $2\cos^2z=1+\cos2z$.
W konsekwencji dla gładkiego rozwiązania
$\partial_t u_3(0)= -\partial_zp_0
=\tfrac18(\cos2x+\cos2y)\sin2z$,
co nie znika tożsamościowo. Trzecia składowa
prędkości nie pozostaje zatem zerowa.

\begin{center}
\begin{tikzpicture}[x=.80cm,y=.80cm,>=Stealth,font=\footnotesize]
\foreach \shift/\az/\zlabel in
 {0/0/{0},5.25/90/{\pi/2},10.5/180/{\pi}} {
  \begin{scope}[shift={(\shift,0)}]
  \fill[gray!5] (0,0) rectangle (4.6,4.6);
  \foreach \i in {0,...,7} {
   \foreach \j in {0,...,7} {
    \pgfmathsetmacro{\xx}{360*(\i+.5)/8}
    \pgfmathsetmacro{\yy}{360*(\j+.5)/8}
    \pgfmathsetmacro{\ww}{2*sin(\xx)*sin(\yy)*cos(\az)}
    \pgfmathsetmacro{\opa}{min(.68,abs(\ww)*.43)}
    \pgfmathsetmacro{\xlo}{4.6*\i/8}
    \pgfmathsetmacro{\ylo}{4.6*\j/8}
    \pgfmathsetmacro{\xhi}{4.6*(\i+1)/8}
    \pgfmathsetmacro{\yhi}{4.6*(\j+1)/8}
    \ifdim\ww pt>0pt
      \fill[red,opacity=\opa] (\xlo,\ylo) rectangle (\xhi,\yhi);
    \else
      \fill[blue,opacity=\opa] (\xlo,\ylo) rectangle (\xhi,\yhi);
    \fi
   }
  }
  \foreach \i in {1,...,8} {
   \foreach \j in {1,...,8} {
    \pgfmathsetmacro{\xx}{360*(\i-.5)/8}
    \pgfmathsetmacro{\yy}{360*(\j-.5)/8}
    \pgfmathsetmacro{\ux}{sin(\xx)*cos(\yy)*cos(\az)}
    \pgfmathsetmacro{\uy}{-cos(\xx)*sin(\yy)*cos(\az)}
    \pgfmathsetmacro{\px}{4.6*(\i-.5)/8}
    \pgfmathsetmacro{\py}{4.6*(\j-.5)/8}
    \draw[-{Stealth[length=2pt]},black!85,line width=.4pt]
      (\px,\py)--++({.36*\ux},{.36*\uy});
   }
  }
  \draw[black!60] (0,0) rectangle (4.6,4.6);
  \node[anchor=south] at (2.3,4.68) {$z=\zlabel$};
  \node[anchor=north] at (0,0) {$0$};
  \node[anchor=north] at (4.6,0) {$2\pi$};
  \end{scope}
}
\node[anchor=west,font=\scriptsize] at (0,-.9)
  {\textcolor{red}{\rule{8pt}{6pt}}\ $\omega_{0,z}>0$\qquad
   \textcolor{blue}{\rule{8pt}{6pt}}\ $\omega_{0,z}<0$\qquad
   strzałki: $(u_{0,1},u_{0,2})$};
\end{tikzpicture}

{\small\emph{Rysunek:} Trzy przekroje pola początkowego wiru
Taylora--Greena, obliczone wprost ze wzorów na $u_0$ i
$\omega_{0,z}$. Strzałki przedstawiają rzuty prędkości na
płaszczyznę $xy$, a czerwień i błękit oznaczają znak pionowej
składowej wirowości (nasycenie barwy rośnie z jej modułem).
Rysunek powstaje w TikZ bez zewnętrznych obrazów i skryptów.}
\end{center}

\begin{center}
\begin{tikzpicture}[x={(.82cm,0cm)},y={(.33cm,.28cm)},
 z={(0cm,1.65cm)},>=Stealth,font=\footnotesize]
 \foreach \h/\az/\lab in {0/0/0,1/90/{\pi/2},2/180/\pi}{
  \fill[gray!5] (0,0,\h)--(6,0,\h)--(6,6,\h)--(0,6,\h)--cycle;
  \foreach \i in {0,...,7}{
   \foreach \j in {0,...,7}{
    \pgfmathsetmacro{\xx}{360*(\i+.5)/8}
    \pgfmathsetmacro{\yy}{360*(\j+.5)/8}
    \pgfmathsetmacro{\ww}{sin(\xx)*sin(\yy)*cos(\az)}
    \pgfmathsetmacro{\opa}{.68*abs(\ww)}
    \pgfmathsetmacro{\lo}{.75*\i}
    \pgfmathsetmacro{\hi}{.75*(\i+1)}
    \pgfmathsetmacro{\bot}{.75*\j}
    \pgfmathsetmacro{\topc}{.75*(\j+1)}
    \ifdim\ww pt>0pt
      \fill[red,opacity=\opa] (\lo,\bot,\h)--(\hi,\bot,\h)
        --(\hi,\topc,\h)--(\lo,\topc,\h)--cycle;
    \else
      \fill[blue,opacity=\opa] (\lo,\bot,\h)--(\hi,\bot,\h)
        --(\hi,\topc,\h)--(\lo,\topc,\h)--cycle;
    \fi
   }
  }
  \foreach \i in {0,...,5}{
   \foreach \j in {0,...,5}{
    \pgfmathsetmacro{\xx}{360*(\i+.5)/6}
    \pgfmathsetmacro{\yy}{360*(\j+.5)/6}
    \pgfmathsetmacro{\ux}{sin(\xx)*cos(\yy)*cos(\az)}
    \pgfmathsetmacro{\uy}{-cos(\xx)*sin(\yy)*cos(\az)}
    \ifnum\az=90\relax
    \else
      \draw[-{Stealth[length=2pt]},black!85,line width=.4pt]
        ({\i+.5},{\j+.5},\h)--++({.42*\ux},{.42*\uy},0);
    \fi
   }
  }
  \draw[black!65] (0,0,\h)--(6,0,\h)--(6,6,\h)--(0,6,\h)--cycle;
  \node[anchor=west] at (6,6,\h) {$z=\lab$};
 }
 \draw[->] (0,0,0)--(6.6,0,0) node[right] {$x$};
 \draw[->] (0,0,0)--(0,6.6,0) node[above] {$y$};
 \draw[->] (0,0,0)--(0,0,2.4) node[above] {$z$};
 \node[below] at (6,0,0) {$2\pi$};
 \node[left] at (0,6,0) {$2\pi$};
\end{tikzpicture}

{\small\emph{Rysunek:} Perspektywiczny widok trzech przekrojów
pola początkowego. Barwy i strzałki wynikają z rzeczywistych
wzorów na $\omega_{0,z}$ i prędkość, a nie ze schematycznych
okręgów. Na $z=\pi/2$ pole zanika; na $z=\pi$ zwroty
prędkości i znaki pionowej składowej wirowości odwracają się
względem $z=0$. Nie jest to wizualizacja ewolucji w czasie.}
\end{center}
\end{przyklad}

\section{Porównanie zagadnień brzegowych i dalsze kierunki}
Na ograniczonym obszarze Lipschitzowskim $\Omega\subset\R^d$
z warunkiem przylegania $u|_{\partial\Omega}=0$
przestrzenie $H,V$ tworzy się przez domknięcie
gładkich bezdywergencyjnych pól o zwartym nośniku
odpowiednio w $L^2(\Omega;\R^d)$ i
$H^1_0(\Omega;\R^d)$ (w drugim przypadku
równoważnie w normie gradientu).
Odzyskiwanie ciśnienia wymaga wówczas analizy warunków
brzegowych, a operator Stokesa nie sprowadza się do
prostego mnożnika Fouriera. Dlatego powyższe dowody
są sformułowane na torusie; ich bezpośrednie przeniesienie
na dowolny obszar bez dodatkowych twierdzeń byłoby nieuzasadnione.

Podstawowa różnica między wymiarami jest następująca:
energia $L^2$ i globalne istnienie rozwiązań Leraya--Hopfa
działają w $d=2,3$; w wymiarze dwa mamy dodatkowo
jednoznaczność i kontrolę wirowości, podczas gdy w wymiarze trzy
pojawia się rozciąganie wirów i warunkowe kryteria regularności.
W kolejnych rozdziałach będziemy badać inne mechanizmy
nieliniowe: swobodne brzegi i homogenizację.

